\documentclass[11pt,a4paper]{article}
\usepackage[T1]{fontenc}
\usepackage{lmodern,amsmath,amssymb,amsthm,booktabs,microtype}
\usepackage[margin=26mm]{geometry}
\usepackage[colorlinks=true,linkcolor=blue,citecolor=blue,urlcolor=blue]{hyperref}
\newtheorem{theorem}{Theorem}[section]
\newtheorem{lemma}[theorem]{Lemma}
\newtheorem{corollary}[theorem]{Corollary}
\newcommand{\Q}{\mathbb Q}\newcommand{\Z}{\mathbb Z}
\DeclareMathOperator{\Res}{Res}\DeclareMathOperator{\Tr}{Tr}
\DeclareMathOperator{\Hom}{Hom}\DeclareMathOperator{\End}{End}
\title{Finiteness of elliptic trace incidence at a fixed slope}
\author{}\date{17 November 2026}
\begin{document}\maketitle
\begin{abstract}
For each fixed input and prepared slope, the first actual detector
trace has only finitely many rational realizations on any fixed
genus-one trace curve. The result allows all independent source
point roles and every rational point on the auxiliary curve. We
compute the rational endomorphism ring of the restriction-of-scalars
surface, prove nonconstancy on every relevant rational abelian
translate using positive detector factors, and apply Mordell--Lang.
No source Mordell--Weil rank assertion is needed. The theorem does
not exclude individual witnesses, a jointly varying slope, or the
untwisted conic criterion.
\end{abstract}

\section{The rational endomorphism ring}
Let
\[
 K=\Q(\sqrt2),\quad d=3-\sqrt2,\quad d'=3+\sqrt2=7/d,
 \qquad E:y^2=x^3-d^2x,\qquad A=\Res_{K/\Q}E.
\]
Write $E^\sigma$ for the conjugate curve.

\begin{lemma}\label{lem:end}
The abelian surface $A$ is simple over $\Q$ and
\[
                         \End_{\Q}(A)=\Z.
\]
In particular $\Hom_{\Q}(A,B)=0$ for every elliptic curve
$B/\Q$.
\end{lemma}
\begin{proof}
Both geometric elliptic endomorphism rings are $\Z[i]$, since
these are characteristic-zero curves with $j=1728$.
Complex conjugation over the real field $K$ fixes precisely $\Z$,
so $\End_K(E)=\End_K(E^\sigma)=\Z$.

The explicit geometric isomorphism
\[
 \phi:E\longrightarrow E^\sigma,\qquad
 (x,y)\longmapsto
       \left(\frac{7x}{d^2},\frac{7\sqrt7\,y}{d^3}\right)
\]
is negated by the automorphism fixing $K(i)$ and negating $\sqrt7$.
This automorphism exists: the quadratic subfields of
$K(i)=\Q(\sqrt2,i)$ are $\Q(\sqrt2)$, $\Q(i)$ and
$\Q(\sqrt{-2})$, so seven is not a square there.
Every geometric homomorphism $E\to E^\sigma$ is
$\phi\circ\alpha$ with $\alpha\in\Z[i]$. The same automorphism
fixes $\alpha$ and negates $\phi$, so a homomorphism defined over
$K$ must be zero. The reverse cross-homomorphism group is also zero.

Over $K$, the restriction of scalars is $E\times E^\sigma$.
Its endomorphism ring over $K$ is therefore $\Z\oplus\Z$.
The nontrivial automorphism of $K/\Q$ exchanges the two diagonal
integers. Descent gives $\End_{\Q}(A)=\Z$. If $A$ were not
$\Q$-simple, Poincar\'e complete reducibility over $\Q$
\cite[Proposition 12.1]{Milne} would give a nontrivial idempotent
in $\End_{\Q}(A)\otimes\Q=\Q$, a contradiction.
A nonzero homomorphism to an elliptic curve would have a
one-dimensional connected kernel over $\Q$, contradicting simplicity.
\end{proof}

\begin{lemma}\label{lem:direction}
Every positive-dimensional abelian subvariety $G\subset A^r$
defined over $\Q$ contains the image of
\[
 A\longrightarrow A^r,\qquad
 Z\longmapsto([a_1]Z,\ldots,[a_r]Z)
\]
for some primitive integer vector $(a_1,\ldots,a_r)$.
\end{lemma}
\begin{proof}
Complete reducibility and simplicity show that $G$ is isogenous
over $\Q$ to a positive power of $A$. A nonzero homomorphism
$A\to G\hookrightarrow A^r$ has integer components by
Lemma~\ref{lem:end}. Dividing their vector by its greatest common
divisor does not change the image, since multiplication by a
nonzero integer on $A$ is surjective. This gives a primitive vector.
\end{proof}

\section{A general incidence theorem}

\begin{theorem}\label{thm:incidence}
Let $B/\Q$ be an elliptic curve, $g\in\Q(B)$ a nonconstant
rational function and $f\in\Q(A^r)$. Suppose $f$ is nonconstant
on every positive-dimensional abelian translate defined over $\Q$
that meets its guarded domain. Then
\[
 \{(R,Q)\in A^r(\Q)\times B(\Q): f(R)=g(Q),
                     \text{ all specified guards hold}\}
\]
is finite.
\end{theorem}
\begin{proof}
Take the reduced closed incidence subvariety $Y\subset A^r\times B$.
Mordell--Weil makes $\Gamma=A^r(\Q)\times B(\Q)$ finitely
generated. Mordell--Lang, in the coset form of
\cite[Theorem 1.1 and equation (1.1)]{GGK}, expresses $Y\cap\Gamma$
as a finite union of cosets of subgroups of $\Gamma$.

Suppose one such coset contains infinitely many guarded points and
choose a guarded rational point $q$ in it. The identity component
$G$ of the Zariski closure of its underlying rational subgroup is a
positive-dimensional abelian subvariety defined over $\Q$, and
$q+G\subset Y$. Thus the argument does not assume that the original
geometric cosets in a Mordell--Lang decomposition were defined over
$\Q$.

Let $G_A$ be the projection of $G$ to $A^r$. If $G$ contains
$0\times B$, then $q+(0\times B)\subset Y$ would make $g$
constant. Otherwise its projection to $G_A$ has finite kernel,
so is an isogeny. Complete reducibility makes $G_A$ isogenous
to a power of $A$; hence its homomorphism space to $B$, and that
of $G$ to $B$, are zero by Lemma~\ref{lem:end}. The second
projection is therefore zero and $G=G_A\times0$. Then $f$ is
constant on the positive-dimensional translate $q_A+G_A$,
contrary to the hypothesis.

Boundary fibers over excluded denominators cause no problem:
$q$ was chosen in the guarded domain, so each translate used in the
contradiction meets the actual incidence. Therefore every relevant
coset contributes only finitely many guarded points.
\end{proof}

The conclusion also holds for a smooth projective genus-one curve
with a rational point, by choosing that point as origin. If such a
curve has no rational point, the incidence is empty.

\section{Positive detector factors on abelian translates}
Fix an input and a slope for which the first detector block is
prepared: its physical weight and comparison scale are nonzero
constants, and every displayed square denominator is positive at
the real places. All source point roles remain independent.

For the conditional source, the three roles are $(A_0,H,U)$ and
\[
\begin{gathered}
 f_0=1/x,\quad z=-x/y,\quad
 r=z(U)/z(H),\quad s=z(U-A_0)/z(H-P),\\
 C_1=f_0(H),\quad C_2=f_0(U),\quad
 C_3=\overline{f_0(H-P)},\quad C_4=\overline{f_0(U-A_0)},\\
 C_5=r-\bar s,\qquad C_6=r+\bar s-2X,\\
 W=\prod_{i=1}^{6}(1+C_i^2/F),\qquad
 V=\Tr_{K/\Q}(W^{-1}),\qquad F>0.
\end{gathered}
\]
The first argument is $t_0=\lambda$, so $F=F(\lambda)$ and
$w=F(\lambda)/(4(\lambda-1)^8)$ are fixed independently of
all three point roles. Every prepared $\lambda>1$ has $F>0$.

For the general source, write the nine independent roles as
\[
 (Q_1,Q_2,Q_3,Q_4,U_0,U_1,U_2,U_3,U_4),\qquad U_0=H.
\]
Its first detector block includes the reciprocal-abscissa tests
\[
 (f_0(U_i))_{i=0}^4,\qquad
 \overline{f_0(U_0-P)},\qquad
 (\overline{f_0(U_i-Q_i)})_{i=1}^4.
\]
Each occurs in a factor $1+f_0(L)^2/(\delta_0^2w_0)$, possibly
with the real slots interchanged. At fixed input and slope,
$a_0=1$ makes $\delta_0,w_0$ constants; $w_0>0$ on the prepared
range. All additional field and rational scalar factors are also
one plus a nonnegative square. These are exactly the first blocks
of the retained quadratic reweighted sources \cite{Reweighted}.

\begin{lemma}\label{lem:nonconstant}
In either source, the first trace $V$ is nonconstant on every
positive-dimensional abelian translate over $\Q$ that meets the
guarded real domain. In particular it is nonconstant at every
prepared positive specialization, with all point roles independent.
\end{lemma}
\begin{proof}
By Lemma~\ref{lem:direction}, such a translate contains a direction
with a primitive integer coefficient vector. In the conditional
case write that vector $(a,h,u)$. If $h$ or $u$ is odd, the
corresponding tested point $H$ or $U$ has odd coefficient. If
both are even, then $a$ is odd and the tested point $U-A_0$ has
odd coefficient $u-a$.

In the general case, if one auxiliary coefficient $u_i$ is odd,
use $U_i$. Otherwise all $u_i$ are even, so a root coefficient
$q_i$ is odd and the tested point $U_i-Q_i$ has odd coefficient.
Thus, in both cases, some reciprocal-abscissa test is evaluated at
\[
                             R+[c]Z
\]
for a fixed real point $R\in A(\mathbb R)$ and an odd integer $c$.

Each real elliptic factor of $A(\mathbb R)$ has two components,
since $x(x^2-d_\sigma^2)$ has three distinct real roots.
Multiplication by an odd integer is surjective on both the identity
tori and the component group. Hence $[c]$ is surjective on
$A(\mathbb R)$. Choose $Z$ tending to a preimage for which both
real slots of $R+[c]Z$ tend to the two-torsion point $(0,0)$.
In both slots the corresponding reciprocal abscissa tends to infinity.

All detector factors are at least one, and the chosen factor gives
$W_\sigma\ge1+c_\sigma f_0(R+[c]Z)_\sigma^2$ with a fixed
$c_\sigma>0$, allowing the slot interchange. Therefore
$V=\sum_\sigma W_\sigma^{-1}\longrightarrow0$.
At the original guarded point, however, $V>0$.
The restricted rational denominator divisors are proper because the
translate contains that guarded point. They can be avoided in
punctured real neighborhoods, so the limiting argument is valid.

The real path tests a rational-function identity; it need not remain
in the square-abscissa or doubled-point component used in source
completeness. Every real open subset of the smooth irreducible
parameter surface $A$ is Zariski dense. Thus the positive value and
zero limit prove nonconstancy on the translate.
\end{proof}

\section{Consequences for the actual trace equations}

\begin{corollary}\label{cor:actual}
Fix the input and a prepared slope in either retained source.
For any fixed smooth genus-one trace curve, only finitely many
rational source point tuples realize the first actual detector
trace on that curve. All three conditional point roles, or all nine
general point roles, may vary independently. The conclusion permits
every rational point on the auxiliary curve and does not assume a
source Mordell--Weil rank or an auxiliary rank bound.
\end{corollary}
\begin{proof}
Apply Theorem~\ref{thm:incidence} with $r=3$ or $9$, $f=V$,
and $g$ the nonconstant trace-coordinate function. Lemma
\ref{lem:nonconstant} gives its translate hypothesis. A rational
point on the genus-one curve may be chosen as origin; if none
exists, there are no incidences.
\end{proof}

For $0<w\le1/4$, the corollary applies to each of
\[
\begin{aligned}
 Y^2&=(V+2)(V^2+w(2-V)),\\
 Y^2&=(V^2-V+2)(V^2+w(2-V)),\\
 Y^2&=(1+w(2-V))(V^2+w(2-V)).
\end{aligned}
\]
These curves are smooth and have rational reference points by
\cite{SmallClasses}. A finite menu of genus-one choices still
gives only finitely many point tuples at each fixed input and slope.
The same conclusion holds for an adaptive coefficient depending
only on that fixed slope, such as $c=\lambda$, whenever the
resulting trace curve is smooth of genus one. Imposing a second
trace equation, the source height equation or additional guards
can only reduce the set.

There can therefore be no infinite fixed-slope propagation of exact
elliptic trace witnesses through independent auxiliary indices,
regardless of the availability of compatible contact congruences or
real approximation. The result does not assert that any fixed-slope
incidence set is empty. It leaves open jointly varying the slope and
point tuple, or making the auxiliary curve itself vary with the point
tuple. The untwisted equation $Y^2=V^2+w(2-V)$ is a conic and is
outside the theorem. No smaller defining-polynomial count is proved.

\begin{thebibliography}{9}\raggedright
\bibitem{Milne}J.~S.~Milne, \emph{Abelian varieties},
Proposition~12.1, author notes.
\url{https://www.jmilne.org/math/xnotes/AVs.pdf}.
\bibitem{GGK}Z.~Gao, T.~Ge and L.~K\"uhne,
\emph{The uniform Mordell--Lang conjecture}, Theorem~1.1 and
equation~(1.1), author manuscript.
\url{https://ziyangjeremygao.github.io/articles/UML.pdf}.
\bibitem{Reweighted}\emph{Quadratic reweightings of native elliptic
detectors}, repository manuscript,
\nolinkurl{papers/pole_parameters/reweighted_elliptic_detectors_20261116_v53.tex}.
\bibitem{SmallClasses}\emph{Small polynomial classes for the actual
detector trace}, repository manuscript,
\nolinkurl{papers/pole_parameters/small_trace_classes_20261117_v58.tex}.
\end{thebibliography}
\end{document}
